%===============================================================================
% LLMWiki Technical Specification - Volume 09: Developer Guide
%===============================================================================
\chapter{Developer Guide}
\label{cha:dev-guide}

This volume is for developers extending LLMWiki: adding parsers, chunkers, embedders, graph backends, index backends, evidence verifiers, planner strategies, tools, and auth providers.

\section{Development Environment}
\label{sec:dev:env}

\subsection{Prerequisites}
\label{sec:dev:prereqs}

\begin{verbatim}
# Rust (for core)
rustup install stable
rustup component add rustfmt clippy

# Python (for ML, parsers, scripts)
python -m venv .venv
source .venv/bin/activate
pip install -r requirements-dev.txt

# Node.js (for Web UI, if applicable)
corepack enable
pnpm install

# System deps
# Ubuntu: sudo apt install libssl-dev libclang-dev pkg-config
# macOS: brew install llvm pkg-config
\end{verbatim}

\subsection{Repository Structure}
\label{sec:dev:structure}

\begin{verbatim}
llmwiki/
-- crates/
   -- llmwiki-core/          # Core types, traits, errors
   -- llmwiki-ingestion/     # Parsers, chunkers, embedders
   -- llmwiki-graph/         # Graph store, entity resolution
   -- llmwiki-index/         # HNSW, BM25, hybrid search
   -- llmwiki-evidence/      # Evidence gate, cross-encoder
   -- llmwiki-planner/       # Query planner, optimizer
   -- llmwiki-runtime/       # Agent runtime, tools, LLM client
   -- llmwiki-api/           # REST/gRPC/WS handlers
   -- llmwiki-cli/           # CLI commands
-- python/
   -- llmwiki/               # Python bindings (pyo3)
   -- scripts/               # Evaluation, benchmarking
-- web/                       # React/Vite admin UI
-- docker/
-- helm/
-- eval/                      # Evaluation datasets, configs
-- docs/                      # This documentation
-- Cargo.toml                 # Workspace root
\end{verbatim}

\subsection{Build and Test}
\label{sec:dev:build-test}

\begin{verbatim}
# Build all (release)
cargo build --release --workspace

# Run unit tests
cargo test --workspace

# Run integration tests (requires Docker)
cargo test --workspace --test integration

# Run specific crate tests
cargo test -p llmwiki-ingestion

# Build Python wheel
maturin develop --release -m python/Cargo.toml

# Lint
cargo fmt --all --check
cargo clippy --workspace -- -D warnings
\end{verbatim}

\section{Extension Points Deep Dive}
\label{sec:dev:extensions}

\subsection{Custom Parser}
\label{sec:dev:parser}

1. Create crate \texttt{llmwiki-parser-myformat}
2. Implement \texttt{Parser} trait:

\begin{verbatim}
// src/lib.rs
use llmwiki_core::{Parser, Artifact, SemanticUnit, ParserSchema, MimeType, Result};
use async_trait::async_trait;

pub struct MyFormatParser;

#[async_trait]
impl Parser for MyFormatParser {
    async fn parse(&self, artifact: Artifact) -> Result<Vec<SemanticUnit>> {
        // 1. Parse artifact.bytes into AST
        // 2. Extract semantic units from AST
        // 3. For each unit: compute content hash, extract metadata
        // 4. Return units
    }

    fn supported_mimes(&self) -> Vec<MimeType> {
        vec!["application/x-myformat".parse().unwrap()]
    }

    fn schema(&self) -> ParserSchema {
        ParserSchema {
            unit_types: vec!["MyUnit".into()],
            metadata_fields: vec!["custom_field".into()],
        }
    }
}

// Register in config
// config.toml:
// [ingestion.parsers.myformat]
// type = "dynamic"
// lib = "llmwiki_parser_myformat"
// entry = "MyFormatParser::new"
\end{verbatim}

3. Add to \texttt{Cargo.toml}:
\begin{verbatim}
[package]
name = "llmwiki-parser-myformat"
cdylib = true  # For dynamic loading

[dependencies]
llmwiki-core = { path = "../../crates/llmwiki-core" }
llmwiki-ingestion = { path = "../../crates/llmwiki-ingestion" }
tree-sitter = "0.22"  # If using tree-sitter
\end{verbatim}

\paragraph{Tree-sitter Grammar Integration.}
If your format has a tree-sitter grammar:
\begin{verbatim}
// build.rs
fn main() {
    tree_sitter_grammar::build("myformat", &["src/grammar.json"]);
}
\end{verbatim}

Then use \texttt{tree-sitter} in parser:
\begin{verbatim}
let mut parser = tree_sitter::Parser::new();
parser.set_language(&tree_sitter_myformat::language())?;
let tree = parser.parse(&source, None)?;
\end{verbatim}

\subsection{Custom Chunker}
\label{sec:dev:chunker}

\begin{verbatim}
use llmwiki_core::{Chunker, SemanticUnit, Chunk, ChunkConfig, Result};

pub struct MyChunker;

impl Chunker for MyChunker {
    fn chunk(&self, unit: SemanticUnit, config: ChunkConfig) -> Vec<Chunk> {
        match config.strategy {
            ChunkStrategy::Hierarchical => self.hierarchical_chunk(unit, config),
            ChunkStrategy::Sliding => self.sliding_chunk(unit, config),
            ChunkStrategy::Semantic => self.semantic_chunk(unit, config),
            _ => vec![Chunk::from_unit(unit)],
        }
    }

    fn config_schema(&self) -> ChunkConfigSchema {
        ChunkConfigSchema {
            fields: vec![
                ConfigField::Int("max_tokens", "Max tokens per chunk"),
                ConfigField::Int("overlap_tokens", "Overlap tokens"),
                ConfigField::Bool("preserve_headers", "Keep markdown headers"),
            ],
        }
    }
}
\end{verbatim}

\subsection{Custom Embedder}
\label{sec:dev:embedder}

\begin{verbatim}
use llmwiki_core::{Embedder, ModelId, Result};
use ndarray::Array2;

pub struct MyEmbedder {
    session: ort::Session,  // ONNX Runtime
    dim: usize,
}

impl MyEmbedder {
    pub fn new(model_path: &str) -> Result<Self> {
        let session = ort::Session::builder()?.commit_from_file(model_path)?;
        let dim = session.outputs[0].dimensions()[1] as usize;
        Ok(Self { session, dim })
    }
}

#[async_trait]
impl Embedder for MyEmbedder {
    async fn embed_batch(&self, texts: &[String]) -> Result<Vec<Vec<f32>>> {
        // Tokenize, pad, run inference, normalize
        let inputs = self.tokenize(texts)?;
        let outputs = self.session.run(inputs)?;
        let embeddings = self.postprocess(outputs)?;
        Ok(embeddings)
    }

    fn dimension(&self) -> usize { self.dim }
    fn model_id(&self) -> ModelId { ModelId::Custom("my-embedder".into()) }
}
\end{verbatim}

\subsection{Custom Graph Store}
\label{sec:dev:graph-store}

Implement \texttt{GraphStore} trait (Vol.~01, \S\ref{iface:graphstore}). Key methods:
\begin{itemize}
  \item \texttt{upsert\_nodes}: Batch insert/update nodes
  \item \texttt{upsert\_edges}: Batch insert/update edges
  \item \texttt{traverse}: Pattern-based traversal
  \item \texttt{neighborhoods}: $k$-hop neighborhoods
\end{itemize}

See \texttt{crates/llmwiki-graph/src/store/rocksdb.rs} for reference implementation.

\subsection{Custom Index Backend}
\label{sec:dev:index-backend}

Implement \texttt{SemanticIndex} trait (Vol.~01, \S\ref{iface:semantic-index}):
\begin{itemize}
  \item \texttt{upsert}: Add/update vectors + metadata
  \item \texttt{search\_vector}: ANN search
  \item \texttt{search\_bm25}: Keyword search
  \item \texttt{search\_hybrid}: Fused search
\end{itemize}

See \texttt{crates/llmwiki-index/src/backend/hnsw\_tantivy.rs}.

\subsection{Custom Evidence Verifier}
\label{sec:dev:evidence-verifier}

Implement \texttt{EvidenceGate} trait (Vol.~01, \S\ref{iface:evidence-gate}):
\begin{verbatim}
pub struct MyEvidenceGate {
    model: CrossEncoder,
    calibrator: Calibrator,
}

impl EvidenceGate for MyEvidenceGate {
    fn verify(&self, query: &str, candidates: Vec<ScoredUnit>, threshold: f32) 
        -> Result<Vec<VerifiedUnit>> {
        // Score each (query, unit) pair
        // Apply calibration
        // Filter by threshold
    }

    fn calibrate(&self, dataset: CalibrationSet) -> Result<CalibrationCurve> {
        // Fit isotonic/Platt/temperature
    }
}
\end{verbatim}

\subsection{Custom Query Planner Strategy}
\label{sec:dev:planner}

Implement \texttt{QueryPlanner} trait (Vol.~01, \S\ref{iface:query-planner}):
\begin{verbatim}
pub struct MyPlanner {
    cost_model: CostModel,
    llm: LLMClient,
}

impl QueryPlanner for MyPlanner {
    fn plan(&self, query: Query, context: PlanContext) -> Result<QueryPlan> {
        // 1. Analyze query intent
        // 2. Decompose into subqueries
        // 3. Choose retrieval strategies per subquery
        // 4. Build DAG with dependencies
        // 5. Optimize order (cost model)
        // 6. Return plan
    }

    fn optimize(&self, plan: QueryPlan, stats: IndexStats) -> Result<QueryPlan> {
        // Reorder, merge, prune based on stats
    }
}
\end{verbatim}

\subsection{Custom Agent Tool}
\label{sec:dev:tool}

\begin{verbatim}
use llmwiki_core::{Tool, ToolInput, ToolOutput, ToolSchema, Result};
use async_trait::async_trait;
use serde::{Deserialize, Serialize};

#[derive(Deserialize)]
struct MyToolInput { query: String, limit: usize }

#[derive(Serialize)]
struct MyToolOutput { results: Vec<String> }

pub struct MyTool;

#[async_trait]
impl Tool for MyTool {
    fn name(&self) -> &str { "my_tool" }
    fn description(&self) -> &str { "Custom tool description" }
    fn schema(&self) -> ToolSchema { ToolSchema::from_type::<MyToolInput>() }
    fn output_schema(&self) -> ToolSchema { ToolSchema::from_type::<MyToolOutput>() }

    async fn execute(&self, input: ToolInput) -> Result<ToolOutput> {
        let params: MyToolInput = serde_json::from_value(input.0)?;
        // Execute logic
        let output = MyToolOutput { results: vec![] };
        Ok(ToolOutput(serde_json::to_value(output)?))
    }
}
\end{verbatim}

Register in config:
\begin{verbatim}
[tools.custom.my_tool]
module = "my_tools"
class = "MyTool"
\end{verbatim}

\subsection{Custom Auth Provider}
\label{sec:dev:auth}

\begin{verbatim}
use llmwiki_core::{AuthProvider, AuthContext, UserInfo, Result};
use async_trait::async_trait;

pub struct MyAuthProvider;

#[async_trait]
impl AuthProvider for MyAuthProvider {
    async fn authenticate(&self, credentials: AuthCredentials) 
        -> Result<AuthContext> {
        // Validate token, API key, etc.
        // Return context with user_id, roles, permissions
    }

    async fn authorize(&self, ctx: &AuthContext, resource: &str, action: &str) 
        -> Result<bool> {
        // Check permissions
    }

    fn extract_credentials(&self, request: &HttpRequest) -> Option<AuthCredentials> {
        // Extract from header, cookie, query param
    }
}
\end{verbatim}

\section{Testing Extensions}
\label{sec:dev:testing}

\subsection{Parser Tests}
\begin{verbatim}
// tests/parser_test.rs
use llmwiki_ingestion::Parser;

#[test]
fn parse_myformat() {
    let parser = MyFormatParser;
    let artifact = Artifact::from_bytes(b"...", "test.myformat");
    let units = parser.parse(artifact).await.unwrap();
    
    assert_eq!(units.len(), 3);
    assert_eq!(units[0].unit_type, "MyUnit");
    assert!(units[0].id.starts_with("sha256:"));
}
\end{verbatim}

\subsection{Integration Tests}
\begin{verbatim}
// tests/integration_test.rs
use llmwiki_core::*;

#[tokio::test]
async fn full_ingestion_query_cycle() {
    let repo = temp_repo();
    llmwiki_init(&repo).await;
    
    // Ingest
    llmwiki_ingest(&repo, "testdata/sample.py").await.unwrap();
    
    // Query
    let result = llmwiki_query(&repo, "What does sample.py do?").await.unwrap();
    
    assert!(!result.answer.is_empty());
    assert!(result.citations.len() > 0);
    assert!(result.citations.iter().all(|c| c.relevance > 0.7));
}
\end{verbatim}

\subsection{Property-Based Testing}
\begin{verbatim}
use proptest::prelude::*;

proptest! {
    #[test]
    fn canonicalization_deterministic(content in "\\PC*") {
        let unit1 = SemanticUnit::new("Test", content.clone());
        unit1.content = content.clone();
        let id1 = unit1.compute_id();
        
        let unit2 = SemanticUnit::new("Test()", content);
        let id2 = unit2.compute_id();
        
        prop_assert_eq!(id1, id2);
    }
}
\end{verbatim}

\section{Debugging and Profiling}
\label{sec:dev:debugging}

\subsection{Logging}
\begin{verbatim}
# Debug parser
RUST_LOG=llmwiki_ingestion=debug llmwiki ingest path/

# Debug query pipeline
RUST_LOG=llmwiki_runtime=debug,llmwiki_planner=trace llmwiki query "test"

# Structured JSON logs
RUST_LOG=info llmwiki server --log-format json
\end{verbatim}

\subsection{Profiling}
\begin{verbatim}
# CPU profiling (perf)
perf record --call-graph dwarf llmwiki ingest large-codebase/
perf report

# Memory profiling
cargo install flamegraph
cargo flamegraph --bin llmwiki -- ingest large-codebase/

# Heap profiling (jemalloc)
MALLOC_CONF=prof:true,prof_prefix:jeprof ./llmwiki ingest ...
jeprof --pdf ./llmwiki jeprof.*.heap > heap.pdf
\end{verbatim}

\section{Contributing}
\label{sec:dev:contributing}

\subsection{Code Style}
\begin{itemize}
  \item Rust: \texttt{rustfmt} (default), \texttt{clippy} (deny warnings)
  \item Python: \texttt{ruff} (format + lint), \texttt{mypy} (strict)
  \item Commits: Conventional Commits (\texttt{feat:}, \texttt{fix:}, \texttt{docs:}, etc.)
\end{itemize}

\subsection{PR Process}
\begin{enumerate}
  \item Fork, create branch \texttt{feat/my-feature}
  \item Write tests for new functionality
  \item Ensure all tests pass: \texttt{cargo test --workspace}
  \item Run evaluation: \texttt{llmwiki eval run --config eval/ci.yaml}
  \item Open PR with description linking to issue
  \item CI must pass (tests, lints, eval regression check)
  \item Maintainer review -> merge
\end{enumerate}

\subsection{Adding a New Volume}
\begin{enumerate}
  \item Create \texttt{docs/Volume\_XX\_Name/Volume\_XX\_Name.tex}
  \item Add \texttt{\textbackslash include\{docs/Volume\_XX\_Name/Volume\_XX\_Name\}} to \texttt{LLMWiki.tex}
  \item Update \texttt{volume-labels.tex} if needed
  \item Run \texttt{make pdf} to verify compilation
\end{enumerate}

\section{Release Process}
\label{sec:dev:release}

\begin{enumerate}
  \item Update \texttt{Cargo.toml} version
  \item Update \texttt{CHANGELOG.md}
  \item Tag: \texttt{git tag v0.X.Y}
  \item GitHub Actions builds and publishes:
    \item Docker images to GHCR
    \item Binaries to GitHub Releases
    \item Python wheels to PyPI
    \item Docs to GitHub Pages
\end{enumerate}

\section{Summary}
\label{sec:dev-guide:summary}

LLMWiki is designed for extensibility at every layer. The trait-based plugin system allows:
\begin{enumerate}
  \item \textbf{Parsers}: New languages/formats without core changes
  \item \textbf{Chunkers}: Domain-specific segmentation strategies
  \item \textbf{Embedders}: Any ONNX-compatible or custom model
  \item \textbf{Graph stores}: RocksDB, SQLite, Kuzu, Neo4j, etc.
  \item \textbf{Indices}: HNSW, FAISS, DiskANN, custom
  \item \textbf{Evidence gates}: Cross-encoder, LLM-judge, heuristic
  \item \textbf{Planners}: Rule-based, LLM-based, hybrid
  \item \textbf{Tools}: Arbitrary computation exposed to planner
  \item \textbf{Auth}: OIDC, API keys, mTLS, custom
\end{enumerate}
All extensions are dynamically loaded, versioned, and testable in isolation.

%===============================================================================
% End of Volume 09
%===============================================================================